% Validation, discussion, limitations, reproducibility and conclusions.

\section{Validation}\label{sec:validation}

\subsection{Analytic checks and oracles}\label{sec:validation-oracles}

A result is believed only when something independent of the engine agrees with it. The first group of checks compares the engine with analytic or manufactured solutions. Total power is conserved by the propagation (Parseval's relation holds for the centred orthonormal transform), and the Fourier form and the point-trap form give the same field at a trap on a bin, on grids of four sizes. The point-trap form matches an explicit evaluation of Eq.~\eqref{eq:delay} with a depth offset and a vortex; the back-projection is the adjoint of the projection, which every weighted update relies on; and the angular spectrum matches an analytic Gaussian beam at one Rayleigh range. For the device, the adjoint of the crosstalk operator passes the inner-product identity on three devices, the kernel matches its definition in Eq.~\eqref{eq:m13}, the blur is the circular convolution with the kernel, the pixel envelope matches a brute-force transform of a pixel-sampled hologram for two fill factors, the crosstalk path reduces to the ideal path when the kernel is a delta, and the efficiency falls monotonically with fewer phase bits, with the first order of a binary grating at the value $4/\pi^2$ within a stated fraction.

The second group are physical limits that no implementation may violate. A random mask reaches an efficiency about the inverse of the trap count (an oracle written separately from the method) and a vortex winds its phase by $2\pi l$. The single-trap reference level of Eq.~\eqref{eq:ceiling} is not among these limits: it is not a bound for many traps (Section~\ref{sec:problem-obj}, Appendix~\ref{app:derivations}) and no test treats it as one. The particle oracle test checks that the simulated bead's power spectrum fits a Lorentzian whose corner frequency is the stiffness of the computed trap divided by $2\pi\gamma$ (within twelve percent) and that the position variance follows equipartition; the dynamics test of the engine checks the same corner frequency within ten percent and that the position histogram follows the Boltzmann law of the computed potential (total variation below 0.06). The oracles live in a module that shares no code with the engine, so that a disagreement is a finding about one of the two sides and no model judges a model. For the gradient methods the closed-form gradient is checked against central finite differences for the powered costs, on shaped targets and on traps, and through the crosstalk operator. The Hungarian solver is checked against brute force on random rectangular matrices, and the incremental field update of the direct search against a full evaluation.

\begin{table}[t]
\centering\footnotesize
\caption{Checks that guard each class of result: the quantity checked, the independent reference it is compared with and the test module that holds it; the tolerance of each is stated in the test and in the design record of the method.}
\label{tab:validation}
\begin{tabular}{@{}>{\raggedright\arraybackslash}p{0.27\textwidth} >{\raggedright\arraybackslash}p{0.33\textwidth} >{\raggedright\arraybackslash}p{0.30\textwidth}@{}}
\toprule
Quantity & Independent reference & Test module\\
\midrule
Fourier form against point-trap form & agreement at on-bin traps on four grid sizes & \texttt{tests/engine/test\_forward}\\
Power conservation & Parseval's relation & \texttt{tests/engine/test\_forward}\\
Depth propagation & analytic Gaussian beam at one Rayleigh range & \texttt{tests/engine/test\_forward}\\
Crosstalk adjoint & inner-product identity, three devices & \texttt{tests/engine/test\_device}\\
Pixel envelope & brute-force transform, two fill factors & \texttt{tests/engine/test\_device}\\
Analytic gradients (M10, M11, M13) & central finite differences & \texttt{tests/methods/optim}\\
Random-mask efficiency & the inverse of the trap count & \texttt{tests/test\_oracles}\\
Vortex winding & $2\pi l$ on a ring of sample points & \texttt{tests/methods/control}\\
Bead spectrum and equipartition & Lorentzian corner frequency, equipartition variance & \texttt{tests/test\_oracles}\\
Bead histogram & Boltzmann law of the computed potential & \texttt{tests/engine/test\_dynamics}\\
Assignment & brute force on random matrices & \texttt{tests/methods/control}\\
Byte-identical artifacts & two bakes of the same inputs & \texttt{tests/test\_manifest}\\
TypeScript engine & fixtures written by the Python engine & \texttt{frontend/src/engine/holo}\\
\bottomrule
\end{tabular}
\end{table}

\subsection{Published values}\label{sec:validation-published}

The second kind of validation reconciles the engine with the published values that the methods replicate (Table~\ref{tab:validation} lists the checks, and Table~\ref{tab:acceptance} the measurements). On the calibration task (Table~\ref{tab:calibration}) part of the ordering of the six baselines of Di Leonardo et al. carries over and part does not. The uniformity of the Gerchberg--Saxton family orders as in the source, plain GS (\calUFour) below the adaptive form (\calUFive) below weighted GS (\calUSix, against the published \pubUSix); the superposition has the lowest uniformity and the random mask the lowest efficiency; and the efficiencies of the three family members are non-increasing in both. The ordering of the two random baselines does not carry over: the published uniformity of random superposition (\pubUTwo) is below that of the random mask (\pubUThree), whereas the engine's smaller grid gives the reverse (\calUTwo{} against \calUThree), which the methods state as the dependence on the pixels per trap. The efficiencies of the Gaussian-beam runs fall below the flat-beam values of the source, and the flat-beam test closes the difference (the single-trap reference level of Eq.~\eqref{eq:ceiling} is the scale of the gap, not a bound). The closed-loop methods are reconciled with the end points of two published curves: the 19 percent to 1.4 percent trap-depth spread of Nogrette et al.~\cite{nogrette} and the 22 percent to 1.59 percent of Kim et al.~\cite{kim2019}, within the factor set in the design document. The device-aware conjugate gradient reproduces the ordering of Schroff et al.~\cite{schroff} and not the published error level; the optimal-transport initialization reproduces the ordering of Swan et al.~\cite{swan} and has an efficiency in the MRAF class and not in the published class; and the random-superposition uniformity and the phase-fixed freeze index are not reproduced, which the corresponding methods state. These differences are kept as results. The case C13 is context for the disordered array of Kim et al. and not a band, because the build measures a lower spread with a Gaussian beam and a camera emulator without read noise, so that the two are not comparable. Where a published number was not re-run, as for the direct search on the quasi-periodic set, it is labelled as quoted.

\subsection{The tests that guard each result}\label{sec:validation-tests}

The repository holds \nTestsTotal{} test functions in the Python test tree, in folders for the engine (\nTests{engine}), the methods (\nTests{methods}), the contracts (\nTests{contracts}), the bake (\nTests{bake}), the learned models (\nTests{learned}), the adversarial plants (\nTests{adversarial}), the notebooks (\nTests{notebooks}), the manuscript checks (\nTests{manuscripts}) and the top-level guards (\nTests{top}), plus the browser-side tests of the live engine. Every acceptance criterion of Table~\ref{tab:ladder} has a test that measures it on the live grid, and the test is the gate named in the criterion; a method whose test fails is not a method of the studio. The adversarial tests plant the failure each gate exists to catch: a fraction written where a percent belongs, a position in metres, a trace over the byte budget, a manifest that differs between two bakes, a committed artifact that changes under a test, a method labelled live that does not fit the budget. Mutation runs confirm that a test can fail: replacing the tuned learning rate of the Wirtinger descent by the source value fails its acceptance test, and replacing the transport seed of the optimal-transport method by a random phase fails its disc and ring tests.

The live engine is held to the Python engine by parity tests over the reference documents that the Python engine writes: \liveFixtures{} documents, each quantity within the tolerance its document states, at most $\liveTol$ absolute, with \liveTests{} browser-side tests in all. A one-iteration shift of the phase freeze of the phase-fixed method, a rule that is easy to get wrong, turns the parity test red, which is how the choice is pinned. The pipeline tests run in a sandbox that fails the test run if a committed artifact changes, and the scripts that regenerate every table and figure of this report read only committed files, so that the numbers in the report are the numbers of the artifacts. The units gate runs in the bake and again over the shipped artifacts; the lane gate and the expected ranges fail the bake on a value outside them. The validation does not establish that the model describes an instrument: that is the limit of Section~\ref{sec:limitations}.

\section{Discussion}\label{sec:discussion}

What the results mean for the problem depends on the target. For point arrays on an ideal device, the integral feedback of the weights and the freeze of the focal phase are the two ingredients that remove the spread: the weighted GS reaches a median spread of \medsig{M06} percent within the budget of the bake, and the phase-fixed form reaches a fraction of a percent when the freeze index is matched to the grid (Table~\ref{tab:bake-sigma}, Figure~\ref{fig:nfix}). The freeze index of the source does not transfer to a smaller SLM, a result the replication produced and the report keeps. The direct search and the Pareto knee reach comparable spreads at a lower efficiency, so the efficiency that the user is willing to pay decides between them, and the front of Section~\ref{sec:results-ablation} shows that on these tasks the price is small because the weighted-GS corner is already near the optimum.

On a real device the picture changes, and it is the main finding of the device cases. The weighted forms equalize the model and not the device, so their spread on the device-limited case stays of the order of ten percent even though their own history falls to a small value; the device-aware weighted GS closes part of that gap on the case and on two geometries of the benchmark, and keeps its start on the lattices of the device because every iterate past the start gives up more than a tenth of the light. The measured support for the method is therefore on a ring and a random set on a device with crosstalk, and not as a replacement for the phase-fixed form on lattices, where the safeguard returns the start. The quantizer and the flicker are the unmodelled part: with them the floor of every method is a few percent, and the straight-through quantizer in the model costs more than it gives on the case that has it.

For shaped potentials, mixed-region freedom trades efficiency for accuracy by design, and the optimal-transport seed gives an error about an order of magnitude lower than the random start at the same efficiency class. For aberrations, the sensorless search recovers an injected set of modes within a small fraction of their amplitude on an ideal camera, and the aberration-robust hologram lowers the 95th percentile of the spread over an ensemble of residuals by a median factor of \accBFourRatio{} at an efficiency cost of \accBFourEffDrop{} percent; against random displacements of the traps its acceptance configuration gives no gain (a ratio of 1.04), while a displacement prior optimized alone for 100 steps reaches 1.48 at 0.10 focal bins at 0.91 of the efficiency, so the budget decides and the curvature of the spot sets only the floor that remains (the companion manuscript of the method gives the measurement). For moving traps, the flicker-bounded sequence keeps the transient intensity higher than linear interpolation in the model of the transient and on most seeds of the longer moves, and not on all; the limit of the seeds that fail is the phase path that the end phases prescribe and not the frame count, so the method needs a phase-path plan before it is a drop-in replacement.

The learned methods are measured on their own held-out sets. The volume network answers in tens of milliseconds and beats the volumetric GS at equal wall-clock on every held-out volume, which makes it useful when the target changes at interaction rate; its advantage is on the accuracy measure of plane images and not on trap uniformity, so a user who needs the spread of Table~\ref{tab:bake-sigma} still runs a weighted method. The warm start is refuted by its control: the early freeze does the work, not the network. That a learned component can be refuted by a control is part of the argument for keeping controls in the protocol. Where each method fails is stated in its section; the user-facing summary is that the studio recommends the phase-fixed form (with its index checked on the user's grid) for point arrays on an ideal device, the device-aware method when the device couples pixels and the geometry is a ring or a random set of the kind that was measured (on the non-lattice quasi-periodic set of the benchmark it is the weakest proposed entry on the C08 device, \gFourBOneSigma{} percent against \gFourPubSigma{} percent for the best published method of that cell, and its support is \benchOnePass{} of \benchOneN{} geometries of the benchmark (\benchOnePassList), one simulated device and \gFourBOneSeeds{} seeds in that cell), the optimal-transport seed for smooth potentials, the sensorless Zernike search before any of them when an aberration is suspected, and the network when the time budget is a few tens of milliseconds.

\section{Limitations}\label{sec:limitations}

DinHot computes and evaluates holograms in simulation: it does not drive an SLM or a camera, and its camera-in-the-loop method runs on an emulated camera whose errors are stated parameters. Device and set-up parameters come from published experiments, so absolute numbers belong to the stated model and not to a measured instrument. The engine is scalar (Fraunhofer and angular spectrum); the high-numerical-aperture case carries a validity flag and the vectorial lane is not built. The fill-factor, four-bit and flicker effects of the device have no published source in this build and are computed only. The source of the crosstalk kernel does not settle whether its width in inverse pixels refers to the SLM pixel or to the computational pixel, one third of it; the engine reads it per SLM pixel, a factor of three away from the other reading, so the kernel is a model of the device and not of a measured one. The transient-intensity model of the flicker-bounded sequence is DinHot's own model of the display between two frames; it ranks methods under one model and is not calibrated against a measured display, and the linear-phase-interpolation baseline is DinHot's implementation of a published description, not the authors' code.

The statistics are limited by the seeds. The unit measurements of Table~\ref{tab:acceptance} use between one and six seeds, the bake three, and the benchmark \benchSeedsMin{} to \benchSeedsMax{} seeds per cell, with a wall-clock bound that lets the count move between runs, so that a count such as the number of geometries on which the front reaches the hypervolume of the published points at 100 iterations is not stable. The uncertainty reported is the observed range or the inter-quartile range, not a confidence interval. The ensemble of the aberration-robust hologram is drawn from the prior it was designed on, so robustness to another distribution is not established. The cgh-net advantage is on a plane-image accuracy measure over generated volumes and the benchmark's volumetric layouts, not on a hardware experiment, and the model is specific to its crop and planes. The learned warm start is refuted and the two-dimensional network ties with GS.

The data are synthetic and the cases few: \nCases{} cases, one of them refused by design, cover the categories but not the diversity of laboratory set-ups. Several of the methods carry reconstructions that are marked in their sections and not confirmed against the original text: the scale convention of the adaptive form, the multiplicative reading of the closed-loop update, the pixel of the kernel width, the step unit of the linear interpolation, the scaling of the front's loss and the accuracy formula of the learned network. The executive figures of the introduction are taken from the cited publications and not re-measured. Finally, the devices and the tweezers of the introduction are neutral-atom arrays and microscopic beads; the optical-trapping literature for cells, neurons and machining is cited for context and no result here is a claim about those uses.

\section{Reproducibility}\label{sec:reproducibility}

The source code, the baked artifacts (manifests, traces, benchmark results, model cards, the ONNX graph and its held-out metrics) and the scripts that regenerate every table and figure of this report from them are in the repository \url{https://github.com/fsantibanezleal/CEFOP_DinHot} (Apache-2.0), under \texttt{manuscripts/holographic-trap-studio/}; the studio runs at \url{https://dinhot.fasl-work.com}. From a clone, with the pipeline environment of the repository, the commands are: the script \texttt{data/acceptance.py} measures the acceptance quantity of each method as its gate computes it; \texttt{data/params.py} reads the keyword defaults of every method; \texttt{data/extract.py} and \texttt{data/extract\_methods.py} write the macro files and the table files from the committed artifacts, the schemas and the catalog; \texttt{figures/make\_figs.py} (with the notebook environment, which has the plotting library) writes the twelve figures as vector PDF; and three passes of \texttt{pdflatex} over \texttt{tex/main.tex} build the report. The canonical bake and the benchmark are explicit release operations (\texttt{python data-pipeline/run.py} and its \texttt{benchmark} command) and take about the times of Section~\ref{sec:protocol-hw}; training is local and writes into the data root, which the repository does not hold. The data sources are the case library of the repository and the committed artifacts; there is no third-party dataset. The version of the product that this report describes is \productVersion{}. The released graph and its card are committed, and the held-out sets are identified by their hash in the model cards.

\section{Conclusions}\label{sec:conclusions}

DinHot compares twenty-four hologram methods on one scalar Fourier model with a device model, \nCases{} cases, a benchmark of \benchRows{} rows and two learned models, and keeps every measured difference from the published sources as a result. The answer to the problem of choosing a method is conditional on the target and the device. Over \hlN{} cases of the bake, among them a device-limited and an aberrated one, the phase-fixed weighted GS leaves a median relative standard deviation of \hlSigmaMseven{} percent (the medians of the single seeds range from \hlSigmaMsevenLo{} to \hlSigmaMsevenHi{}) against \hlSigmaMfour{} percent (\hlSigmaMfourLo{} to \hlSigmaMfourHi{}) for plain GS, means over three seeds, but only with a freeze index matched to the grid; on a device with crosstalk the weighted forms equalize the model and not the device, and the device-aware weighted GS lowers the spread by a factor of \bakeCeightRatioThree{} on the ratio of the means of the device-limited case. Of the five proposed methods, three meet their criteria within a measured scope, one meets a narrowed criterion and one is refuted; the volume network beats its baseline at equal wall-clock on all \cghN{} held-out volumes with a median ratio of \cghMed{}. The next steps are the measured hardware loop with a real SLM and camera, the vectorial focus, a calibrated device model, and a label-based or rule-changing design for the warm start, none of which is in this build.

\begin{thebibliography}{99}\small\raggedright
\bibitem{manetsch} H.~J.~Manetsch, G.~Nomura, E.~Bataille, X.~Lv, K.~H.~Leung and M.~Endres, ``A tweezer array with 6,100 highly coherent atomic qubits,'' \emph{Nature} 647, 60--67, 2025. \url{https://doi.org/10.1038/s41586-025-09641-4}.
\bibitem{ultraprecise} Y.~T.~Chew \emph{et al.}, ``Ultraprecise holographic optical tweezer array,'' \emph{Phys. Rev. A} 110(5), 053518, 2024. \url{https://doi.org/10.1103/PhysRevA.110.053518}.
\bibitem{continuous3000} N.-C.~Chiu \emph{et al.}, ``Continuous operation of a coherent 3,000-qubit system,'' \emph{Nature} 646(8087), 1075--1080, 2025. \url{https://doi.org/10.1038/s41586-025-09596-6}.
\bibitem{logical} D.~Bluvstein \emph{et al.}, ``Logical quantum processor based on reconfigurable atom arrays,'' \emph{Nature} 626, 58--65, 2024. \url{https://doi.org/10.1038/s41586-023-06927-3}.
\bibitem{airearr} R.~Lin \emph{et al.}, ``AI-enabled parallel assembly of thousands of defect-free neutral atom arrays,'' \emph{Phys. Rev. Lett.} 135(6), 060602, 2025. \url{https://doi.org/10.1103/2ym8-vs82}.
\bibitem{feedbackeq} S.~Wang, Y.~Hu, Y.~Zhou, P.~Lan and Z.~Xu, ``Feedback intensity equalization algorithm for multi-spots holographic tweezer,'' arXiv:2407.17049, 2024 (preprint; no journal version found). \url{https://arxiv.org/abs/2407.17049}.
\bibitem{slmcal} P.~Schroff, E.~Haller, S.~Kuhr and A.~La Rooij, ``Rapid stochastic spatial light modulator calibration and pixel crosstalk optimization,'' \emph{Opt. Express} 32(27), 48957, 2024. \url{https://doi.org/10.1364/OE.539548}.
\bibitem{tweezerbio} C.~J.~Bustamante, Y.~R.~Chemla, S.~Liu and M.~D.~Wang, ``Optical tweezers in single-molecule biophysics,'' \emph{Nat. Rev. Methods Primers} 1, 25, 2021. \url{https://doi.org/10.1038/s43586-021-00021-6}.
\bibitem{neuralens} A.~R.~Mardinly \emph{et al.}, ``Precise multimodal optical control of neural ensemble activity,'' \emph{Nat. Neurosci.} 21(6), 881--893, 2018. \url{https://doi.org/10.1038/s41593-018-0139-8}.
\bibitem{machining} Z.~Kuang \emph{et al.}, ``High throughput diffractive multi-beam femtosecond laser processing using a spatial light modulator,'' \emph{Appl. Surf. Sci.} 255(5), 2284--2289, 2008. \url{https://doi.org/10.1016/j.apsusc.2008.07.091}.
\bibitem{neuralholo} Y.~Peng, S.~Choi, N.~Padmanaban and G.~Wetzstein, ``Neural holography with camera-in-the-loop training,'' \emph{ACM Trans. Graph.} 39(6), 1--14, 2020. \url{https://doi.org/10.1145/3414685.3417802}.
\bibitem{dileonardo} R.~Di Leonardo, F.~Ianni and G.~Ruocco, ``Computer generation of optimal holograms for optical trap arrays,'' \emph{Opt. Express} 15(4), 1913, 2007. \url{https://doi.org/10.1364/OE.15.001913}.
\bibitem{bianchi2010} S.~Bianchi and R.~Di Leonardo, ``Real-time optical micro-manipulation using optimized holograms generated on the GPU,'' \emph{Comput. Phys. Commun.} 181(8), 1444--1448, 2010. \url{https://doi.org/10.1016/j.cpc.2010.04.012}.
\bibitem{montes2006} M.~Montes-Usategui, E.~Pleguezuelos, J.~Andilla and E.~Mart\'in-Badosa, ``Fast generation of holographic optical tweezers by random mask encoding of Fourier components,'' \emph{Opt. Express} 14(6), 2101, 2006. \url{https://doi.org/10.1364/OE.14.002101}.
\bibitem{polin} M.~Polin, K.~Ladavac, S.-H.~Lee, Y.~Roichman and D.~G.~Grier, ``Optimized holographic optical traps,'' \emph{Opt. Express} 13(15), 5831, 2005. \url{https://doi.org/10.1364/OPEX.13.005831}.
\bibitem{kim2019} D.~Kim \emph{et al.}, ``Large-scale uniform optical focus array generation with a phase spatial light modulator,'' \emph{Opt. Lett.} 44(12), 3178, 2019. \url{https://doi.org/10.1364/OL.44.003178}.
\bibitem{ebadi2021} S.~Ebadi \emph{et al.}, ``Quantum phases of matter on a 256-atom programmable quantum simulator,'' \emph{Nature} 595, 227--232, 2021. \url{https://doi.org/10.1038/s41586-021-03582-4}.
\bibitem{pasienski} M.~Pasienski and B.~DeMarco, ``A high-accuracy algorithm for designing arbitrary holographic atom traps,'' \emph{Opt. Express} 16(3), 2176, 2008. \url{https://doi.org/10.1364/OE.16.002176}.
\bibitem{harte} T.~Harte, G.~D.~Bruce, J.~Keeling and D.~Cassettari, ``Conjugate gradient minimisation approach to generating holographic traps for ultracold atoms,'' \emph{Opt. Express} 22(22), 26548, 2014. \url{https://doi.org/10.1364/OE.22.026548}.
\bibitem{chakra} P.~Chakravarthula, Y.~Peng, J.~Kollin, H.~Fuchs and F.~Heide, ``Wirtinger holography for near-eye displays,'' \emph{ACM Trans. Graph.} 38(6), 1--13, 2019. \url{https://doi.org/10.1145/3355089.3356539}.
\bibitem{kingma} D.~P.~Kingma and J.~Ba, ``Adam: a method for stochastic optimization,'' 3rd International Conference on Learning Representations, 2015; arXiv:1412.6980. \url{https://arxiv.org/abs/1412.6980}.
\bibitem{nogrette} F.~Nogrette \emph{et al.}, ``Single-atom trapping in holographic 2D arrays of microtraps with arbitrary geometries,'' \emph{Phys. Rev. X} 4(2), 021034, 2014. \url{https://doi.org/10.1103/PhysRevX.4.021034}.
\bibitem{schroff} P.~Schroff, A.~La Rooij, E.~Haller and S.~Kuhr, ``Accurate holographic light potentials using pixel crosstalk modelling,'' \emph{Sci. Rep.} 13, 3252, 2023. \url{https://doi.org/10.1038/s41598-023-30296-6}.
\bibitem{persson} M.~Persson, D.~Engstr\"om and M.~Goks\"or, ``Reducing the effect of pixel crosstalk in phase only spatial light modulators,'' \emph{Opt. Express} 20(20), 22334, 2012. \url{https://doi.org/10.1364/OE.20.022334}.
\bibitem{noll} R.~J.~Noll, ``Zernike polynomials and atmospheric turbulence,'' \emph{J. Opt. Soc. Am.} 66(3), 207, 1976. \url{https://doi.org/10.1364/JOSA.66.000207}.
\bibitem{whyte} G.~Whyte and J.~Courtial, ``Experimental demonstration of holographic three-dimensional light shaping using a Gerchberg--Saxton algorithm,'' \emph{New J. Phys.} 7, 117, 2005. \url{https://doi.org/10.1088/1367-2630/7/1/117}.
\bibitem{lee2017} W.~Lee, H.~Kim and J.~Ahn, ``Defect-free atomic array formation using the Hungarian matching algorithm,'' \emph{Phys. Rev. A} 95(5), 053424, 2017. \url{https://doi.org/10.1103/PhysRevA.95.053424}.
\bibitem{lpi} I.~H.~A.~Knottnerus, Y.~C.~Tseng, A.~Urech, R.~J.~C.~Spreeuw and F.~Schreck, ``Parallel assembly of neutral atom arrays with an SLM using linear phase interpolation,'' \emph{SciPost Phys.} 19(4), 118, 2025. \url{https://doi.org/10.21468/SciPostPhys.19.4.118}.
\bibitem{eybposh} M.~H.~Eybposh, N.~W.~Caira, M.~Atisa, P.~Chakravarthula and N.~C.~P\'egard, ``DeepCGH: 3D computer-generated holography using deep learning,'' \emph{Opt. Express} 28(18), 26636, 2020. \url{https://doi.org/10.1364/OE.399624}.
\bibitem{ren2024} Z.~Ren \emph{et al.}, ``Creation of a tweezer array for cold atoms utilizing a generative neural network,'' \emph{APL Quantum} 1(4), 046111, 2024. \url{https://doi.org/10.1063/5.0219170}.
\bibitem{shi2021} L.~Shi, B.~Li, C.~Kim, P.~Kellnhofer and W.~Matusik, ``Towards real-time photorealistic 3D holography with deep neural networks,'' \emph{Nature} 591, 234--239, 2021. \url{https://doi.org/10.1038/s41586-020-03152-0}.
\bibitem{swan} H.~Swan, A.~Torchylo, M.~J.~Van de Graaff, J.~Rudolph and J.~M.~Hogan, ``High-fidelity holographic beam shaping with optimal transport and phase diversity,'' \emph{Opt. Express} 33(3), 6290--6303, 2025. \url{https://doi.org/10.1364/OE.540901}.
\bibitem{torchylo} A.~Torchylo, H.~Swan, L.~Tellez and J.~M.~Hogan, ``A fast, large-scale optimal transport algorithm for holographic beam shaping,'' \emph{Opt. Express} 34(4), 6796--6803, 2026. \url{https://doi.org/10.1364/OE.588337}.
\bibitem{machu} Y.~Machu, G.~Creutzer, C.~Sayrin and M.~Brune, ``Full-field-of-view aberration correction for large arrays of focused beams,'' \emph{Opt. Express}, 2026. \url{https://doi.org/10.1364/OE.595662}; preprint arXiv:2512.15967.
\bibitem{grimm} R.~Grimm, M.~Weidem\"uller and Y.~B.~Ovchinnikov, ``Optical dipole traps for neutral atoms,'' \emph{Adv. At. Mol. Opt. Phys.} 42, 95--170, 2000. \url{https://doi.org/10.1016/S1049-250X(08)60186-X}.
\bibitem{holland} N.~Holland, D.~Stuart, O.~Barter and A.~Kuhn, ``Benchmarking modern algorithms to holographically create optical tweezers for laser-cooled atoms,'' \emph{J. Mod. Opt.} 65(18), 2133--2141, 2018. \url{https://doi.org/10.1080/09500340.2018.1499978}.
\bibitem{mahajan1983} V.~N.~Mahajan, ``Strehl ratio for primary aberrations in terms of their aberration variance,'' \emph{J. Opt. Soc. Am.} 73(6), 860--861, 1983. \url{https://doi.org/10.1364/JOSA.73.000860}.
\bibitem{harada} Y.~Harada and T.~Asakura, ``Radiation forces on a dielectric sphere in the Rayleigh scattering regime,'' \emph{Opt. Commun.} 124(5--6), 529--541, 1996. \url{https://doi.org/10.1016/0030-4018(95)00753-9}.
\bibitem{bouloumis} T.~D.~Bouloumis and S.~Nic Chormaic, ``From far-field to near-field micro- and nanoparticle optical trapping,'' \emph{Appl. Sci.} 10(4), 1375, 2020. \url{https://doi.org/10.3390/app10041375}.
\bibitem{tolic} S.~F.~Toli\'c-N{\o}rrelykke, E.~Sch\"affer, J.~Howard, F.~S.~Pavone, F.~J\"ulicher and H.~Flyvbjerg, ``Calibration of optical tweezers with positional detection in the back focal plane,'' \emph{Rev. Sci. Instrum.} 77(10), 103101, 2006. \url{https://doi.org/10.1063/1.2356852}.
\bibitem{norrelykke} S.~F.~N{\o}rrelykke and H.~Flyvbjerg, ``Power spectrum analysis with least-squares fitting: amplitude bias and its elimination, with application to optical tweezers and atomic force microscope cantilevers,'' \emph{Rev. Sci. Instrum.} 81(7), 075103, 2010. \url{https://doi.org/10.1063/1.3455217}.
\bibitem{blackman} D.~Blackman and S.~Vigna, ``Scrambled linear pseudorandom number generators,'' \emph{ACM Trans. Math. Softw.} 47(4), 1--32, 2021. \url{https://doi.org/10.1145/3460772}.
\bibitem{lizana} A.~Lizana \emph{et al.}, ``Time fluctuations of the phase modulation in a liquid crystal on silicon display: characterization and effects in diffractive optics,'' \emph{Opt. Express} 16(21), 16711, 2008. \url{https://doi.org/10.1364/OE.16.016711}.
\bibitem{moreno} I.~Moreno, B.~K.~Gutierrez, M.~M.~S\'anchez-L\'opez, J.~A.~Davis, H.~P.~Khanal and D.~M.~Cottrell, ``Diffraction efficiency of stepped gratings using high phase-modulation spatial light modulators,'' \emph{Opt. Lasers Eng.} 126, 105910, 2020. \url{https://doi.org/10.1016/j.optlaseng.2019.105910}.
\bibitem{companion} F.~Santib\'a\~nez-Leal, ``Device-aware weighted Gerchberg--Saxton and flicker-bounded phase-continuous sequences for holographic optical trap arrays,'' Zenodo, concept DOI \url{https://doi.org/10.5281/zenodo.23788733}.
\end{thebibliography}

\appendix

\section{Contract fields}\label{app:fields}

The tables of Section~\ref{sec:data} list the top-level fields of each contract. This appendix lists the nested fields of the larger contracts, generated from the schema files by the same script: the path of the field (a dot descends into an object, brackets into an array item, and braces into a map keyed by an identifier such as a method id), its type, whether the schema requires it, its unit where the name carries one, and the range or constant that the schema states. The schema files are the single source of the contracts: the Python producers and consumers validate against them in their tests, the TypeScript types are generated from them, and a change to a field bumps the version of its contract by the semantic-version rule. Fields that the schema leaves unconstrained show a dash in the range column; a dash in the unit column marks an identifier, a count or a dimensionless quantity.

Table~\ref{tab:fields-params-record} lists the parameter record, one case as the ingestion accepts it, and Table~\ref{tab:fields-case-manifest} the case manifest, the authoritative record of one baked case. Table~\ref{tab:fields-trace} gives the trace, the replay artifact of one case; Table~\ref{tab:fields-benchmark-results} the benchmark results, the protocol's table by method, geometry and budget; Table~\ref{tab:fields-model-metrics} the held-out metrics of the model slots; and Table~\ref{tab:fields-readiness} the readiness record, which states what is real in the build and which the app reads.

\par\medskip\noindent\textbf{Contract \texttt{params-record}}\par\nopagebreak
\begingroup\scriptsize
\begin{longtable}{@{}l l c l l@{}}
\toprule Field & Type & Req. & Unit & Range\\\midrule\endhead
\texttt{case\_id} & string & yes & - & pattern\\
\texttt{setup} & object & yes & - & -\\
\texttt{setup.wavelength\_nm} & number & yes & nm & 300 to 2000\\
\texttt{setup.focal\_mm} & number & yes & mm & 0.5 to 1000\\
\texttt{setup.na} & number & yes & - & 0.05 to 1.5\\
\texttt{setup.n\_imm} & number & yes & - & 1.0 to 1.6\\
\texttt{setup.n\_pix} & enum & yes & - & \{64, 128, 256, 512, 1024\}\\
\texttt{setup.pitch\_um} & number & yes & $\mu$m & 1 to 50\\
\texttt{setup.pad} & integer & yes & - & 1 to 4\\
\texttt{setup.beam} & object & yes & - & -\\
\texttt{device} & object or null & yes & - & -\\
\texttt{target} & object & yes & - & -\\
\texttt{aberration} & object or null & yes & - & -\\
\texttt{methods} & array of string & yes & - & 1 to 24 items\\
\texttt{reference\_method} & string & yes & - & pattern\\
\texttt{budget} & object & yes & - & -\\
\texttt{budget.iterations} & integer & yes & - & 1 to 5000\\
\texttt{budget.time\_s} & number or null & yes & s & 0.01 to 3600\\
\texttt{seed} & integer & yes & - & 0 to 2147483647\\
\texttt{feedback} & object or null & yes & - & -\\
\texttt{dynamics} & object or null & yes & - & -\\
\texttt{rearrangement} & object or null & yes & - & -\\
\texttt{expect} & object & yes & - & -\\
\bottomrule\end{longtable}
\endgroup

\par\medskip\noindent\textbf{Contract \texttt{case-manifest}}\par\nopagebreak
\begingroup\scriptsize
\begin{longtable}{@{}l l c l l@{}}
\toprule Field & Type & Req. & Unit & Range\\\midrule\endhead
\texttt{schema} & const & yes & - & = dinhot.manifest/v1\\
\texttt{case\_id} & string & yes & - & length >= 1\\
\texttt{title} & object & yes & - & -\\
\texttt{title.en} & string & yes & - & length >= 1\\
\texttt{title.es} & string & yes & - & length >= 1\\
\texttt{category} & object & yes & - & -\\
\texttt{category.en} & string & yes & - & length >= 1\\
\texttt{category.es} & string & yes & - & length >= 1\\
\texttt{real\_or\_synthetic} & enum & yes & - & \{real, synthetic\}\\
\texttt{expected\_band} & object & yes & - & -\\
\texttt{expected\_band.en} & string & yes & - & length >= 1\\
\texttt{expected\_band.es} & string & yes & - & length >= 1\\
\texttt{engine} & object & yes & - & -\\
\texttt{engine.package} & const & yes & - & = dinhot-holo\\
\texttt{engine.version} & string & yes & - & pattern\\
\texttt{engine.model} & const & yes & - & = scalar-fourier\\
\texttt{params} & object & yes & - & -\\
\texttt{params.setup} & object & yes & - & -\\
\texttt{params.setup.wavelength\_nm} & number & yes & nm & 300 to 2000\\
\texttt{params.setup.focal\_mm} & number & yes & mm & 0.5 to 1000\\
\texttt{params.setup.na} & number & yes & - & 0.05 to 1.5\\
\texttt{params.setup.n\_imm} & number & yes & - & 1.0 to 1.6\\
\texttt{params.setup.n\_pix} & enum & yes & - & \{64, 128, 256, 512, 1024\}\\
\texttt{params.setup.pitch\_um} & number & yes & $\mu$m & 1 to 50\\
\texttt{params.setup.pad} & integer & yes & - & 1 to 4\\
\texttt{params.setup.beam} & object & yes & - & -\\
\texttt{params.device} & object or null & yes & - & -\\
\texttt{params.target} & object & yes & - & -\\
\texttt{params.aberration} & object or null & yes & - & -\\
\texttt{params.methods} & array of string & yes & - & 1 to 24 items\\
\texttt{params.reference\_method} & string & yes & - & pattern\\
\texttt{params.budget} & object & yes & - & -\\
\texttt{params.budget.iterations} & integer & yes & - & 1 to 5000\\
\texttt{params.budget.time\_s} & number or null & yes & s & 0.01 to 3600\\
\texttt{params.seed} & integer & yes & - & 0 to 2147483647\\
\texttt{params.feedback} & object or null & yes & - & -\\
\texttt{params.dynamics} & object or null & yes & - & -\\
\texttt{params.rearrangement} & object or null & yes & - & -\\
\texttt{params.expect} & object & yes & - & -\\
\texttt{seed} & integer & yes & - & -\\
\texttt{artifact} & object & yes & - & -\\
\texttt{artifact.path} & string & yes & - & pattern\\
\texttt{artifact.format} & const & yes & - & = json\\
\texttt{artifact.trace\_schema} & const & yes & - & = dinhot.trace/v1\\
\texttt{artifact.bytes} & integer & yes & - & 1 to\\
\texttt{lane} & enum & yes & - & \{live, precompute\}\\
\texttt{gate} & object & yes & - & -\\
\texttt{gate.lane} & enum & yes & - & \{live, precompute\}\\
\texttt{gate.pure\_python} & boolean & yes & - & -\\
\texttt{gate.wheels} & array of string & yes & - & -\\
\texttt{gate.trace\_bytes} & integer & yes & bytes & 1 to\\
\texttt{gate.run\_ms\_budget} & number & yes & - & >0 to\\
\texttt{gate.trace\_bytes\_budget} & integer & yes & - & 1 to\\
\texttt{gate.reasons} & array of string & yes & - & -\\
\texttt{flags} & array of object & yes & - & -\\
\texttt{methods} & object & yes & - & -\\
\texttt{methods\{id\}.lane} & enum & yes & - & \{live, baked\}\\
\texttt{methods\{id\}.efficiency} & number & yes & fraction & -\\
\texttt{methods\{id\}.uniformity} & number & yes & fraction & -\\
\texttt{methods\{id\}.rel\_std\_pct} & number & yes & percent & -\\
\texttt{methods\{id\}.zero\_order} & number & yes & fraction & -\\
\texttt{methods\{id\}.ghost} & number & yes & fraction & -\\
\texttt{methods\{id\}.iterations} & integer & yes & - & -\\
\texttt{methods\{id\}.converged} & boolean & yes & - & -\\
\texttt{models} & object & yes & - & -\\
\texttt{expect} & object & yes & - & -\\
\texttt{anchor} & object & yes & - & -\\
\texttt{anchor.source} & object & yes & - & -\\
\texttt{anchor.source.en} & string & yes & - & length >= 1\\
\texttt{anchor.source.es} & string & yes & - & length >= 1\\
\texttt{anchor.url} & string or null & yes & - & pattern\\
\bottomrule\end{longtable}
\endgroup

\par\medskip\noindent\textbf{Contract \texttt{trace}}\par\nopagebreak
\begingroup\scriptsize
\begin{longtable}{@{}l l c l l@{}}
\toprule Field & Type & Req. & Unit & Range\\\midrule\endhead
\texttt{schema} & const & yes & - & = dinhot.trace/v1\\
\texttt{case\_id} & string & yes & - & length >= 1\\
\texttt{grid} & object & yes & - & -\\
\texttt{grid.n\_pix} & integer & yes & - & -\\
\texttt{grid.pad} & integer & yes & - & -\\
\texttt{grid.bin\_um} & number & yes & $\mu$m & -\\
\texttt{grid.half\_field\_um} & number & yes & $\mu$m & -\\
\texttt{traps} & array of object & yes & - & 0 to 4096 items\\
\texttt{traps[].id} & integer & yes & - & 0 to 4095\\
\texttt{traps[].x\_um} & number & yes & $\mu$m & -100000.0 to 100000.0\\
\texttt{traps[].y\_um} & number & yes & $\mu$m & -100000.0 to 100000.0\\
\texttt{traps[].z\_um} & number & yes & $\mu$m & -100000.0 to 100000.0\\
\texttt{traps[].weight} & number & yes & fraction & 0.01 to 100\\
\texttt{traps[].charge} & integer & yes & - & -20 to 20\\
\texttt{shape} & object or null & yes & - & -\\
\texttt{methods} & object & yes & - & -\\
\texttt{methods\{id\}.history} & array of object & yes & - & 0 to 200 items\\
\texttt{methods\{id\}.history[].iteration} & integer & yes & - & -\\
\texttt{methods\{id\}.history[].efficiency} & number & yes & fraction & -\\
\texttt{methods\{id\}.history[].uniformity} & number & yes & fraction & -\\
\texttt{methods\{id\}.history[].rel\_std\_pct} & number & yes & percent & -\\
\texttt{methods\{id\}.final} & object & yes & - & -\\
\texttt{methods\{id\}.final.efficiency} & number & yes & fraction & -\\
\texttt{methods\{id\}.final.uniformity} & number & yes & fraction & -\\
\texttt{methods\{id\}.final.rel\_std\_pct} & number & yes & percent & -\\
\texttt{methods\{id\}.final.zero\_order} & number & yes & fraction & -\\
\texttt{methods\{id\}.final.ghost} & number & yes & fraction & -\\
\texttt{methods\{id\}.trap\_intensity} & array of number & yes & - & -\\
\texttt{reference} & object & yes & - & -\\
\texttt{reference.method} & string & yes & - & pattern\\
\texttt{reference.phase\_u8} & string & yes & - & pattern\\
\texttt{reference.focal} & object & yes & - & -\\
\texttt{reference.focal.size} & integer & yes & - & 1 to 192\\
\texttt{reference.focal.bin\_um} & number & yes & $\mu$m & -\\
\texttt{reference.focal.db\_floor} & number & yes & - & -\\
\texttt{reference.focal.intensity\_db\_u8} & string & yes & - & pattern\\
\texttt{reference.axial} & object or null & yes & - & -\\
\texttt{feedback} & object or null & yes & - & -\\
\texttt{dynamics} & object or null & yes & - & -\\
\texttt{rearrangement} & object or null & yes & - & -\\
\texttt{summary} & object & yes & - & -\\
\texttt{summary.best\_method} & string & yes & - & pattern\\
\texttt{summary.notes} & object & yes & - & -\\
\texttt{summary.notes.en} & string & yes & - & length >= 1\\
\texttt{summary.notes.es} & string & yes & - & length >= 1\\
\bottomrule\end{longtable}
\endgroup

\par\medskip\noindent\textbf{Contract \texttt{benchmark-results}}\par\nopagebreak
\begingroup\scriptsize
\begin{longtable}{@{}l l c l l@{}}
\toprule Field & Type & Req. & Unit & Range\\\midrule\endhead
\texttt{schema} & const & yes & - & = dinhot.benchmark/v1\\
\texttt{engine\_version} & string & yes & - & pattern\\
\texttt{protocol} & object & yes & - & -\\
\texttt{protocol.grids} & array of integer & yes & - & -\\
\texttt{protocol.seeds} & integer & yes & - & -\\
\texttt{protocol.budgets} & object & yes & - & -\\
\texttt{protocol.budgets.iterations} & array of integer & yes & - & -\\
\texttt{protocol.budgets.time\_s} & array of number & yes & s & -\\
\texttt{protocol.device} & object & yes & - & -\\
\texttt{protocol.device.en} & string & yes & - & length >= 1\\
\texttt{protocol.device.es} & string & yes & - & length >= 1\\
\texttt{geometries} & array of object & yes & - & -\\
\texttt{geometries[].id} & string & yes & - & pattern\\
\texttt{geometries[].name} & object & yes & - & -\\
\texttt{geometries[].name.en} & string & yes & - & length >= 1\\
\texttt{geometries[].name.es} & string & yes & - & length >= 1\\
\texttt{geometries[].n\_traps} & integer & yes & - & -\\
\texttt{geometries[].dims} & enum & yes & - & \{2, 3\}\\
\texttt{rows} & array of object & yes & - & -\\
\texttt{rows[].method} & string & yes & - & pattern\\
\texttt{rows[].geometry} & string & yes & - & pattern\\
\texttt{rows[].device} & enum & yes & - & \{ideal, device\}\\
\texttt{rows[].budget\_kind} & enum & yes & - & \{iterations, time, converged\}\\
\texttt{rows[].budget} & number & yes & - & -\\
\texttt{rows[].params} & enum & yes & - & \{default, tuned\}\\
\texttt{rows[].n\_seeds} & integer & yes & - & -\\
\texttt{rows[].efficiency} & object & yes & fraction & -\\
\texttt{rows[].efficiency.median} & number & yes & - & -\\
\texttt{rows[].efficiency.q1} & number & yes & - & -\\
\texttt{rows[].efficiency.q3} & number & yes & - & -\\
\texttt{rows[].uniformity} & object & yes & fraction & -\\
\texttt{rows[].uniformity.median} & number & yes & - & -\\
\texttt{rows[].uniformity.q1} & number & yes & - & -\\
\texttt{rows[].uniformity.q3} & number & yes & - & -\\
\texttt{rows[].rel\_std\_pct} & object & yes & percent & -\\
\texttt{rows[].rel\_std\_pct.median} & number & yes & - & -\\
\texttt{rows[].rel\_std\_pct.q1} & number & yes & - & -\\
\texttt{rows[].rel\_std\_pct.q3} & number & yes & - & -\\
\texttt{rows[].time\_s} & object & yes & s & -\\
\texttt{rows[].time\_s.median} & number & yes & - & -\\
\texttt{rows[].time\_s.q1} & number & yes & - & -\\
\texttt{rows[].time\_s.q3} & number & yes & - & -\\
\texttt{rows[].iterations} & object & yes & - & -\\
\texttt{rows[].iterations.median} & number & yes & - & -\\
\texttt{rows[].iterations.q1} & number & yes & - & -\\
\texttt{rows[].iterations.q3} & number & yes & - & -\\
\texttt{pareto} & array of object & yes & - & -\\
\texttt{pareto[].geometry} & string & yes & - & pattern\\
\texttt{pareto[].method} & const & yes & - & = B5\\
\texttt{pareto[].points} & array of object & yes & - & -\\
\texttt{pareto[].points[].lambda} & number & yes & - & -\\
\texttt{pareto[].points[].efficiency} & number & yes & fraction & -\\
\texttt{pareto[].points[].uniformity} & number & yes & fraction & -\\
\texttt{skipped} & array of object & no & - & -\\
\texttt{skipped[].method} & string & yes & - & pattern\\
\texttt{skipped[].geometry} & string & yes & - & pattern\\
\texttt{skipped[].reason} & object & yes & - & -\\
\texttt{skipped[].reason.en} & string & yes & - & length >= 1\\
\texttt{skipped[].reason.es} & string & yes & - & length >= 1\\
\bottomrule\end{longtable}
\endgroup

\par\medskip\noindent\textbf{Contract \texttt{model-metrics}}\par\nopagebreak
\begingroup\scriptsize
\begin{longtable}{@{}l l c l l@{}}
\toprule Field & Type & Req. & Unit & Range\\\midrule\endhead
\texttt{schema} & const & yes & - & = dinhot.model-metrics/v1\\
\texttt{engine\_version} & string & yes & - & pattern\\
\texttt{slots} & object & yes & - & -\\
\texttt{slots.cgh-net} & object & no & - & -\\
\texttt{slots.cgh-net.held\_out\_n} & integer & yes & - & -\\
\texttt{slots.cgh-net.accuracy\_vs\_gs} & object & yes & - & -\\
\texttt{slots.cgh-net.time\_ms} & object & yes & ms & -\\
\texttt{slots.cgh-net.baseline} & const & yes & - & = M19\\
\texttt{slots.warmstart-model} & object & no & - & -\\
\texttt{slots.warmstart-model.held\_out\_n} & integer & yes & - & -\\
\texttt{slots.warmstart-model.iterations\_to\_1pct} & object & yes & - & -\\
\texttt{slots.warmstart-model.final\_rel\_std\_pct} & object & yes & percent & -\\
\bottomrule\end{longtable}
\endgroup

\par\medskip\noindent\textbf{Contract \texttt{readiness}}\par\nopagebreak
\begingroup\scriptsize
\begin{longtable}{@{}l l c l l@{}}
\toprule Field & Type & Req. & Unit & Range\\\midrule\endhead
\texttt{schema} & const & yes & - & = example.readiness/v1\\
\texttt{product\_version} & string & yes & - & pattern\\
\texttt{catalog\_sha256} & string & yes & - & pattern\\
\texttt{contracts} & array of object & yes & - & 1 to items\\
\texttt{contracts[].id} & string & yes & - & pattern\\
\texttt{contracts[].version} & string & yes & - & pattern\\
\texttt{contracts[].status} & enum & yes & - & \{real, active\}\\
\texttt{contracts[].provenance} & enum & yes & - & \{real, stub, simulator, fixture, baseline, ...\}\\
\texttt{contracts[].state} & enum & yes & - & \{available, partial, stale, provider, training, ...\}\\
\texttt{contracts[].reason} & object or null & yes & - & -\\
\texttt{contracts[].missing} & array of string & yes & - & -\\
\bottomrule\end{longtable}
\endgroup


\section{The window, the envelope and the efficiency ceiling}\label{app:derivations}

Three relations of Section~\ref{sec:problem} are used without a full derivation in the body. The first is the unaliased window, derived for this product and not quoted from a paper. The SLM samples the phase on a grid of pitch $p$, so that its field repeats in the focal plane with the period $\lambda f/p$ of the grating equation; a centred grid of that period spans from $-X_{1/2}$ to $X_{1/2}$ with $X_{1/2}=\lambda f/(2p)$, and a trap at or beyond the bound folds back onto the grid and cannot be told from a trap inside it, which is why the first contract refuses it. One bin of the transform padded by $\mathrm{pad}$ is the period divided by the padded size, $\Delta x=\lambda f/(\mathrm{pad}\,n_{\mathrm{pix}}\,p)$.

The second is the pixel-aperture envelope. A square pixel of width $a=\sqrt{\mathrm{fill}}\,p$ is a rectangular aperture, whose Fourier transform is a cardinal sine of argument $a x/(\lambda f)$ along each axis. At the edge of the window the argument is $\sqrt{\mathrm{fill}}/2$, so that the intensity is
\begin{equation}
\mathcal E(X_{1/2})=\operatorname{sinc}^2\!\Big(\frac{\sqrt{\mathrm{fill}}}{2}\Big),
\label{eq:envelope-edge}
\end{equation}
about 0.4 for a fill factor of one, the value that the pixel-aperture test confirms against a brute-force transform of a pixel-sampled hologram. The envelope has unit value at the centre, and the power lost to dead area is not applied, so the efficiencies of the report are fractions of the power that reaches the focal plane.

The third is the efficiency ceiling, which is a bound for one trap only. The efficiency of Eq.~\eqref{eq:metrics} is the sum of the trap intensities, each a fraction of the input power. A hologram with the best possible phase at a single trap gives the field $V=\frac1{N_{\mathrm{pix}}}\sum_jA_j$, because all the terms of Eq.~\eqref{eq:field} are then in phase; its intensity is the square of the mean amplitude, and no phase pattern can do better at that trap. For a flat beam $A_j=1$ and the bound is one; for a Gaussian beam it is the square of the mean amplitude of Eq.~\eqref{eq:ceiling}, below one. For many traps the only bound for every beam and every phase is $e\le1$, and $\overline A^2$ is exceeded: a numerical maximization for a Gaussian beam with $\overline A^2=0.81$ reaches $0.995$ at ten traps. The tests use $\overline A^2$ as a reference level (the Gaussian tests state 0.9 of it), so those thresholds are relative to the single-trap limit and not to a proven ceiling of the array.
